\section{Infrastrukturkrav}
\label{sec:krav-infrastruktur}

Infrastrukturkravene beskriver hvordan mellomvaren skal pakkes, 
konfigureres og kjøres i et test- eller driftsmiljø.

\subsection{Containerisering (I1.1)}
\label{subsec:krav-i11}

Krav I1.1 innebærer at mellomvaren skal kunne kjøres som et 
Docker-image. Formålet er å gjøre løsningen enklere å distribuere, 
teste og kjøre uten manuell installasjon av rammeverk og avhengigheter 
på vertsmaskinen.

Kravet vurderes som oppfylt dersom applikasjonen kan bygges som et 
Docker-image, startes som container og konfigureres uten at imaget må 
bygges på nytt. Oppsettet skal også dokumenteres slik at HDO kan kjøre 
løsningen videre etter prosjektperioden.

\subsection{Stateless design og skalerbarhet (I1.2)}
\label{subsec:krav-i12}

Krav I1.2 innebærer at mellomvaren skal være stateless. I denne 
oppgaven tolkes dette som at VideoAPI ikke skal lagre rom-, bruker- 
eller sesjonstilstand lokalt per instans. Dersom slik tilstand er 
nødvendig, skal den ligge i en delt ekstern komponent eller hos de 
eksterne videomotorene.

Kravet innebærer også at miljøavhengig konfigurasjon, som backend-URLer, 
API-nøkler og token-konfigurasjon, skal kunne settes gjennom 
miljøvariabler.

Kravet evalueres gjennom arkitekturgjennomgang, kontroll av 
konfigurasjonshåndtering og testing av containerisert kjøring. 
Begrensninger knyttet til horisontal skalering diskuteres i 
kapittel~\ref{chap:diskusjon}.